\section{第13章：热力学基础（13.5练习题）}
\sourceproblem{选择1}{235}{241}
\problemtext 活塞快速外移导致气体膨胀，判断过程性质。
\workings 气体内部一般不能维持均匀平衡压强，非准静态，不能用平衡态组成的单一 $p$-$V$ 曲线表示，选C。
\crosscheck 可绘制外压或平均压强历程，但不等于平衡态路径。
\sourceproblem{选择2}{235}{241}
\problemtext 同初态、同末体积的等压、等温、绝热膨胀，比较吸热量。
\workings 等压路径压强最高，$W_p>W_T$，且 $\Delta U_p>0,\Delta U_T=0$；绝热 $Q=0$。故等压吸热最多，选A。
\crosscheck 按图示准静态路径与正热容模型。
\sourceproblem{选择3}{235}{241}
\problemtext 同初态三种过程使体积加倍，比较温变绝对值。
\workings 等压 $T_2=2T_1$，等温不变；可逆绝热 $T_2=T_1 2^{1-\gamma}$，故 $0<|\Delta T_{\rm ad}|<T_1$。等压最大、等温最小，选D。
\crosscheck 取定热容、gamma大于1的理想气体。
\sourceproblem{选择4}{235}{241}
\problemtext 理想气体绝热自由膨胀的终态。
\workings $Q=W=0$，由第一定律 $\Delta U=0$，故T不变而V增大，p减小，选A。
\crosscheck 初末均为平衡态；自由膨胀不是可逆绝热过程。
\sourceproblem{选择5}{235}{241}
\problemtext 选择Q、内能增量及对外功均为负的过程。
\workings 等压压缩 $W=p\Delta V<0$，$T$ 降低，$\Delta U<0$，$Q=\Delta U+W<0$，选D。
\crosscheck 等体降压功为零，绝热膨胀热量为零。
\sourceproblem{选择6}{235}{241}
\problemtext 图示a至b等体升压，再至c等压膨胀，判断三个能量量的符号。
\workings $p_cV_c>p_aV_a$，故 $\Delta U>0$；第二段体积增加，$W>0$，因此 $Q>0$，选B。
\crosscheck 根据原图两段箭头，而非仅由末态判断功。
\sourceproblem{选择7}{236}{242}
\problemtext 系统吸热后内能是否必增加。
\workings $\Delta U=Q-W$，对外功可小于、等于或大于Q，内能可增、恒定或减，选D。
\crosscheck 吸热量不是状态函数。
\sourceproblem{选择8}{236}{242}
\problemtext 等物质的量同初态氦、氮、水蒸气定容等量吸热。
\workings 刚性模型 $C_{V,m}=iR/2$，三种气体自由度3、5、6，$\Delta T=Q/(\nu C_{V,m})$ 不同；$\Delta p=\nu R\Delta T/V$ 亦不同，选C。
\crosscheck 水为非线性多原子分子，忽略振动。
\sourceproblem{选择9}{236}{242}
\problemtext 相同刚性容器内氢、氦同初态，氦吸热3 J，求氢同温升吸热。
\workings 等物质的量且定容，热容量比 $5:3$，氢需5 J，选C。
\crosscheck 经典刚性模型。
\sourceproblem{选择10}{236}{242}
\problemtext 等容同初态氨、氢，氢吸热5 J，求氨同温升吸热。
\workings 氨为非线性刚性分子，i=6，氢i=5，故 $Q_{\rm NH_3}=5(6/5)=6$ J，选A。
\crosscheck 不是按原子数直接取自由度。
\sourceproblem{选择11}{236}{242}
\problemtext 1 mol单原子理想气体已知始末p、V、T而未知路径及种类，求可确定量。
\workings $\Delta U=3R(T_B-T_A)/2$ 可确定，选B；W、Q依赖路径，质量还需摩尔质量。
\crosscheck 内能变化为状态量之差。
\sourceproblem{选择12}{236}{242}
\problemtext 直线a(4 atm,2 L)至b(2 atm,4 L)，判断功热。
\workings 初末pV均8 atm L，故 $\Delta U=0$；体积增大，$W>0$，$Q=W>0$，选B。
\crosscheck 虽初末温度相同，直线路径不一定等温。
\sourceproblem{选择13}{236}{242}
\problemtext 定体摩尔热容视为常数时内能变化公式适用范围。
\workings 理想气体U仅依赖T，$\Delta U=\nu C_{V,m}\Delta T$ 适用于任何初末平衡态的过程，选D。
\crosscheck 路径可以非准静态；变热容时用积分。
\sourceproblem{选择14}{236}{242}
\problemtext 图示acb吸热500 J，求闭合acbda净吸热（图续243页）。
\workings 图中a、b的pV相同，故 $W_{acb}=500$ J。b至d等体，d至a等压压缩，$W_{bda}=-4\times10^5(3\times10^{-3})=-1200$ J；循环 $\Delta U=0$，故 $Q_{\rm net}=-700$ J，选B。
\crosscheck 图中压强、体积轴分别有10的5次方及负3次方因子。
\sourceproblem{选择15}{237}{243}
\problemtext 刚性双原子理想气体等压膨胀W/Q。
\workings $W=\nu R\Delta T$，$Q=7\nu R\Delta T/2$，比为2/7，选D。
\crosscheck 取常温刚性五自由度。
\sourceproblem{选择16}{237}{243}
\problemtext 绝热膨胀至2V、等体恢复初温、等温压缩回原态的循环。
\workings 绝热膨胀路径位于初温等温线下，膨胀功小于压缩功绝对值，故 $W_{\rm net}<0$，$Q_{\rm net}=W_{\rm net}<0$，选A。
\crosscheck 循环内能变化为零。
\sourceproblem{选择17}{237}{243}
\problemtext 正循环的p-V图特征。
\workings 顺时针闭合曲线的 $\oint p\,dV>0$，选C。
\crosscheck 按通常横V、纵p的坐标与简单循环。
\sourceproblem{选择18}{237}{243}
\problemtext 图示V-T循环AB、BC、CA中哪个吸热。
\workings AB沿过原点直线且V、T均增，等压吸热；BC等体降温放热，CA等温压缩放热，选A。
\crosscheck 图中纵轴V、横轴T，不是p-V图。
\sourceproblem{选择19}{237}{243}
\problemtext 37摄氏度及27摄氏度热源间热机的最大效率。
\workings $\eta_{\max}=1-300/310=1/31\simeq3.23\%$。原书选项25%、50%、75%、91.74%均不符。
\crosscheck 须用绝对温度；不能直接以摄氏温度比计算，不强配选项。
\sourceproblem{选择20}{237}{243}
\problemtext 400 K与300 K热源间声称吸热1800 J、放热800 J、做功1000 J的卡诺机。
\workings 能量守恒成立，但 $\eta=1000/1800=5/9>1-300/400=1/4$，不可能，选D。
\crosscheck 第一定律成立并不足以保证过程可实现。
\sourceproblem{选择21}{237}{243}
\problemtext 分别提高高温热源或降低低温热源相同温差，比较卡诺效率增量。
\workings $\Delta\eta_1=T_2\Delta T/[T_1(T_1+\Delta T)]$，$\Delta\eta_2=\Delta T/T_1$，前者较小，选B。
\crosscheck 要求两种调整后的温度均物理可行。
\sourceproblem{选择22}{237}{243}
\problemtext 绝热膨胀至2V0、等压恢复初温、等温压缩回初态（选项续244页）。
\workings 绝热段及随后的低压等压膨胀均在初温等温线下方，净功为负；循环 $\Delta U=0$，故净吸热为负，选A。
\crosscheck 末等温压缩从更大体积回初体积。
\sourceproblem{选择23}{238}{244}
\problemtext 卡诺循环两绝热线下阴影面积S1、S2比较。
\workings 两绝热段连接相同的高低温等温线，$Q=0$，功的绝对值均等于 $|\Delta U|$，故 $S_1=S_2$，选B。
\crosscheck 面积指图示绝热路径功的绝对值。
\sourceproblem{选择24}{238}{244}
\problemtext 四幅等体、等温、等压、绝热路径组成的设想循环中判断可行者。
\workings 图B先等体升温、再绝热膨胀、最后沿较低温等温线压缩可闭合；图A同起点膨胀绝热线不应在等温线上方；图C、D两条不同可逆绝热线不能相交。选B。
\crosscheck 对固定量理想气体的可逆绝热线，pV的gamma次幂为常量；不同常量曲线不交。
\sourceproblem{选择25}{238}{244}
\problemtext 按热力学第二定律辨别热功转换、传热与自发过程说法。
\workings 宏观自发非平衡过程不可逆，选D。A遗漏循环且无其他影响条件；B遗漏不借助外界作用条件；C把不可逆误为绝对不能反向。
\crosscheck 可逆极限过程及微观涨落不等同于通常宏观自发过程。
\sourceproblem{选择26}{238}{244}
\problemtext 辨别自由膨胀与热功转换的第二定律陈述。
\workings 气体可自发自由膨胀而不自发宏观收缩，选C。
\crosscheck 热泵可借助外功使热由低温向高温传递，等温膨胀可使吸热全变为功但有体积改变。
\sourceproblem{选择27}{238}{244}
\problemtext 单一热源接触的理想气体等温膨胀，吸热全用于做功是否违反定律。
\workings $\Delta U=0,Q=W$，第一定律成立；气体体积改变，不是单热源循环且无其他影响，不违反第二定律，选C。
\crosscheck 不能把非循环膨胀当成循环热机。
\sourceproblem{选择28}{238}{244}
\problemtext 比较四种关于热机效率、卡诺效率及定律的陈述。
\workings 第一定律不能单独排除效率1，也不能单独证明任意工质卡诺效率；第二定律排除单热源全热变功循环。对理想气体可借第一定律、物态方程及可逆绝热关系计算卡诺效率。按教材语境乙、丁对，选D。
\crosscheck 丁并非仅凭抽象能量守恒即可成立，还依赖理想气体模型及可逆路径条件。
\sourceproblem{选择29}{239}{245}
\problemtext 判断热功转换及热传递四项陈述。
\workings 在通常宏观热机及有限温差自发传热语境，(2)热机效率小于1、(4)高温向低温自发传热不可逆正确，选A。(1)(3)遗漏限制条件。
\crosscheck 可逆等温传热为零温差极限，不是有限温差自发传热。
\sourceproblem{选择30}{239}{245}
\problemtext 图示内能E随V沿过原点的直线上升，判断过程。
\workings 按理想气体热运动内能零点及定热容模型，$U=\nu C_{V,m}T$，$U/V=(C_{V,m}/R)p$。图示比值恒定，等压，选A。
\crosscheck 若内能含任意加性常数，则需先扣除该常数；原图按教材零点。
\sourceproblem{填空1}{239}{245}
\problemtext p-V图中平衡态、准静态过程及准静态循环的表示。
\workings 分别为一点、一条连续曲线、一条闭合曲线。
\crosscheck 采用固定物质的量的简单可压缩系统。
\sourceproblem{填空2}{239}{245}
\problemtext 热平衡理想气体不随时间变化的三个宏观量及不断变化的微观量。
\workings 压强、体积、温度；分子位置、速度及瞬时动能等。
\crosscheck 宏观量取平衡统计平均，微小涨落并非绝对为零。
\sourceproblem{填空3}{239}{245}
\problemtext 图中上路径1、下路径2，二者间面积S1，下路径以下面积S2，求a-1-b及a-2-b-1-a的功。
\workings $W_{a1b}=S_1+S_2$，$W_{\rm cycle}=S_2-(S_1+S_2)=-S_1$。
\crosscheck 循环方向逆时针。
\sourceproblem{填空4}{239}{245}
\problemtext 同初态等体吸热416 J，等压至相同末温吸热582 J，求等压功。
\workings 等体 $\Delta U=416$ J，同末温故等压内能增量亦416 J，$W=582-416=166$ J。
\crosscheck 不要求气体具体种类。
\sourceproblem{填空5}{239}{245}
\problemtext 绝热容器内搅拌液体升温，判断热量、外界做功、内能增加与零的关系。
\workings $Q=0$，外界对系统做功 $W_{\rm on}>0$，$\Delta U=W_{\rm on}>0$（初末宏观动能不变）。
\crosscheck 搅拌机械功不是热量。
\sourceproblem{填空6}{239}{245}
\problemtext 内能增加的方式及其改变量依赖因素。
\workings 可由外界做功或吸热实现；内能改变只决定于初末状态，与具体路径无关。
\crosscheck 吸热与外界做功可同时发生。
\sourceproblem{填空7}{239}{245}
\problemtext 等温压缩外界做功绝对值W1，再绝热膨胀返回初体积对外做功绝对值W2，求总热和内能增量。
\workings 第一段 $\Delta U_1=0,Q_1=-|W_1|$；第二段 $Q_2=0,\Delta U_2=-|W_2|$。总 $Q=-|W_1|,\Delta U=-|W_2|$。
\crosscheck 返回原体积不代表回到原温度。
\sourceproblem{填空8}{240}{246}
\problemtext 定压摩尔热容大于定体摩尔热容的原因。
\workings 等压升温时除增加内能还对外做膨胀功，$C_{p,m}-C_{V,m}=R$。
\crosscheck 理想气体、仅体积功。
\sourceproblem{填空9}{240}{246}
\problemtext 同初态三条膨胀曲线AB等压、AC等温、AD绝热及内能变化。
\workings 取 $p_A,V_1$，三曲线为 $p=p_A$、$p=p_AV_1/V$、$p=p_A(V_1/V)^\gamma$，均从A起，至V2时B最高、C居中、D最低。AB内能增加，AD减少，AC不变。
\crosscheck 此解析式可直接用于绘制原题要求曲线；gamma大于1。
\sourceproblem{填空10}{240}{246}
\problemtext 同初态膨胀至2V1，等压、等温、绝热中比较功、内能增量、吸热。
\workings 等压路径压强最高、温度上升，故三者均为等压过程最大。
\crosscheck 绝热吸热为零，等温内能变化为零。
\sourceproblem{填空11}{240}{246}
\problemtext 理想气体定容、定温、定压吸热的能量去向。
\workings 定容：全部增加内能；定温：全部对外做功；定压：一部分增加内能，一部分对外做功。
\crosscheck 假定固定气体、仅体积功、热容为正。
\sourceproblem{填空12}{240}{246}
\problemtext 图示p-T循环1至2曲线、2至3竖直、3至1水平，判热量符号。
\workings 1至2的温度升高且压强降低，体积 $V=\nu RT/p$ 增大，故 $\Delta U>0,W>0,Q>0$。2至3等温升压压缩，放热；3至1等压降温压缩，放热。
\crosscheck 依原图曲线从1向2降压升温，不额外假设其为某种等值过程。
\sourceproblem{填空13}{240}{246}
\problemtext 三图等压线、等温线、绝热线中比较a、b内能。
\workings 图1 b体积较大，同压温度较高，b内能大；图2 a、b同温，内能相等；图3 a体积较小，可逆绝热 $TV^{\gamma-1}$ 恒定，a内能大。
\crosscheck 按原图a左b右。
\sourceproblem{填空14}{240}{246}
\problemtext 10 mol单原子气体压缩时外界做功209 J，升温1 K，求内能增量和吸热。
\workings $\Delta U=3(10)R/2=124.65$ J；$W=-209$ J，$Q=124.65-209=-84.35$ J。
\crosscheck 负吸热即放热84.35 J。
\sourceproblem{填空15}{240}{246}
\problemtext 等压对外做功200 J，单原子及刚性双原子分别需吸热多少。
\workings $Q=(C_{p,m}/R)W$，单原子500 J，双原子700 J。
\crosscheck 采用3、5自由度。
\sourceproblem{填空16}{241}{247}
\problemtext 1 mol刚性双原子气体等压膨胀做功W，求温变及吸热。
\workings $\Delta T=W/R$，$Q_p=7W/2$。
\crosscheck 仅有体积功。
\sourceproblem{填空17}{241}{247}
\problemtext 文字给A(2p1,V1)、B(2p1,V2)，图却标B(p1,2V1)，求直线路径功及内能变化。
\workings 按图 $W=(2p_1+p_1)(2V_1-V_1)/2=3p_1V_1/2$，端点pV相同，$\Delta U=0$。若严格按文字，B压强同A，则 $W=2p_1(V_2-V_1)$，定热容时 $\Delta U=(C_{V,m}/R)2p_1(V_2-V_1)$。
\crosscheck 原页文字和图不一致，保留两套答案；未给气体种类，文字解不强定热容。
\sourceproblem{填空18}{241}{247}
\problemtext 1 mol单原子气体由(p1,V1)沿图中直线至(p2,V2)，求功及吸热。
\workings $W=(p_1+p_2)(V_2-V_1)/2$，$\Delta U=3(p_2V_2-p_1V_1)/2$，$Q=W+\Delta U$。
\crosscheck 直线延长通过原点，尚有p1/V1=p2/V2；上述式亦适用于任意直线。
\sourceproblem{填空19}{241}{247}
\problemtext 常温刚性分子气体等压过程，求W/Q及内能增量/Q。
\workings 以i为二次自由度数，$W/Q=2/(i+2)$，$\Delta U/Q=i/(i+2)$。
\crosscheck 题未给分子种类，不能指定唯一数值。
\sourceproblem{填空20}{241}{247}
\problemtext 1 mol气体等压升温1 K，内能增加20.78 J，求功及吸热。
\workings $W=R\Delta T=8.31$ J，$Q=20.78+8.31=29.09$ J。
\crosscheck 直接使用题给内能值。
\sourceproblem{填空21}{241}{247}
\problemtext 3 mol气体500 K等温压强从6 atm降至3 atm，求吸热。
\workings $Q=W=3(8.31)(500)\ln2\simeq8640$ J。
\crosscheck 按准静态等温过程，若不可逆则仅初末态不足以确定功热。
\sourceproblem{填空22}{241}{247}
\problemtext 氢、氦同p、V、T，比较质量、内能、分别等压吸等量热时功。
\workings 等物质的量，质量比1:2，内能比5:3；$W/Q=R/C_{p,m}$，故功比 $(2/7):(2/5)=5:7$。
\crosscheck 刚性氢分子、单原子氦模型。
\sourceproblem{填空23}{241}{247}
\problemtext 低温热源27摄氏度，卡诺效率40%，提高至50%时高温源温度增量。
\workings $T_h=300/(1-0.4)=500$ K；新高温源600 K，需升100 K。
\crosscheck 固定低温热源。
\sourceproblem{填空24}{241}{247}
\problemtext 290 g空气卡诺机，热源27与负73摄氏度，高温等温膨胀体积比2.718，求效率及净功。
\workings $\nu=10$ mol，$T_h=300,T_c=200$ K，$\eta=1/3$，$W=\nu R(T_h-T_c)\ln2.718\simeq8.309\times10^3$ J。
\crosscheck 2.718近似自然常数e；非恰等于e。
\sourceproblem{填空25}{241}{247}
\problemtext 热源727与527摄氏度，最大效率热机每循环吸热2000 J，求功。
\workings $T_h=1000,T_c=800$ K，$\eta_{\max}=0.2$，$W=400$ J。
\crosscheck 使用绝对温度。
\sourceproblem{填空26}{241}{247}
\problemtext 图示顺时针矩形循环，p从10至40 Pa，V从1至4立方厘米，求净吸热。
\workings $Q_{\rm net}=W_{\rm net}=(40-10)(4-1)10^{-6}=9.0\times10^{-5}$ J。
\crosscheck 原图V单位为立方厘米而非立方米。
\sourceproblem{填空27}{242}{248}
\problemtext 1 mol气体T-V图AB水平、BC竖直、CA可逆绝热，A(T1,V1)、B(T2,V2)，求C状态。
\workings 图示 $T_1=T_2$，$V_C=V_2$。由CA绝热 $T_CV_2^{\gamma-1}=T_1V_1^{\gamma-1}$，$T_C=T_1(V_1/V_2)^{\gamma-1}$，$p_C=RT_C/V_2$。
\crosscheck 沿用图示与定gamma模型。
\sourceproblem{填空28}{242}{248}
\problemtext 标准态理想气体可逆绝热膨胀至3V0，求末温和压强。
\workings $T=T_0 3^{1-\gamma}$，$p=p_0 3^{-\gamma}$。
\crosscheck 题目称绝热，教材默认准静态可逆；不可逆绝热不能只由体积比确定。
\sourceproblem{填空29}{242}{248}
\problemtext 热力学第二定律克劳修斯及开尔文表述。
\workings 不可能把热量从低温物体传到高温物体而不引起其他影响；不可能从单一热源吸热，使之全部变为功而不引起其他影响。
\crosscheck 两个表述均含其他影响的限制；开尔文表述适用于循环热机。
\sourceproblem{填空30}{242}{248}
\problemtext 两种第二定律表述分别说明哪类过程不可逆。
\workings 教材所指为功耗散成热的过程及有限温差下热量自发由高温向低温传递的过程。
\crosscheck 并非所有理想化热力学过程均不可逆；可逆过程是理想极限。
\sourceproblem{计算1}{242}{248}
\problemtext 刚性双原子及单原子气体等压膨胀分别做功2 J，求吸热。
\workings $Q=C_{p,m}W/R$，双原子 $Q=(7/2)2=7$ J，单原子 $Q=(5/2)2=5$ J。
\crosscheck 固定量理想气体，定热容模型。
\sourceproblem{计算2}{242}{248}
\problemtext 双原子气体准静态路径pV平方为常数a，从V1膨胀至V2，求功、内能变化、吸热。
\workings $p=a/V^2$，故 $W=a(1/V_1-1/V_2)>0$。刚性双原子 $U=5pV/2$，所以 $\Delta U=(5a/2)(1/V_2-1/V_1)=-5W/2$；$Q=-3W/2$。
\crosscheck 膨胀同时降温并放热；若振动激发需改用实际热容积分。
\sourceproblem{计算3}{242}{248}
\problemtext 图示直线a(3 atm,1 L)经b(2 atm,2 L)至c(1 atm,3 L)，求功、内能增量、吸热。
\workings $W=(3+1)(3-1)/2=4$ atm L，即 $405.2$ J。初末pV均3 atm L，故 $\Delta U=0$，$Q=405.2$ J。
\crosscheck 1 atm L等于101.3 J；中途b的pV为4 atm L，整个过程并非等温。

\sourceproblem{计算4}{243}{249}
\problemtext 单原子理想气体由A等压膨胀至B，再绝热膨胀至C。原图A、B压强40万Pa，C压强10万Pa，体积分别2、3.49、8立方厘米，求总功和内能变化。
\workings 图示 $p_A=p_B=4\times10^5$ Pa，$p_C=10^5$ Pa。$W_{AB}=p_A(V_B-V_A)=0.596$ J。$W_{BC}=-\Delta U_{BC}=\frac32(p_BV_B-p_CV_C)=0.894$ J。故 $W=1.490$ J，$\Delta U=\frac32(p_CV_C-p_AV_A)=0$。
\crosscheck 体积轴为cm³。B点体积3.49为舍入值，可逆绝热关系给约3.482 cm³；总功按题图读数。
\sourceproblem{计算5}{243}{249}
\problemtext 刚性绝热容器分两室，各有1 mol氦和氧气，由无摩擦导热活塞隔开，初温400、600 K且初压相同为p0，求平衡温压。
\workings 整体无外界功热交换，$\frac32R(400)+\frac52R(600)=4RT$，故 $T=525$ K。总容积 $V=R(400+600)/p_0$ 不变；末态共同压强 $p=2RT/V=1.05p_0=1.06365\times10^5$ Pa。
\crosscheck 氧气按刚性双原子模型，活塞热容量及最终宏观动能忽略；两室内部功、热在整体能量账中相消。
\sourceproblem{计算6}{243}{249}
\problemtext 单原子理想气体三角循环，A、B、C的(V,p)按图为(1,1)、(2,3)、(2,1)，单位分别为10的负3次方立方米、10的5次方Pa，求各段及循环能量。
\workings 采用 $W=\int p\,dV,\ \Delta U=3\Delta(pV)/2,\ Q=\Delta U+W$。AB直线段 $W=200$ J、$\Delta U=750$ J、$Q=950$ J；BC等容 $W=0$、$\Delta U=-600$ J、$Q=-600$ J；CA等压 $W=-100$ J、$\Delta U=-150$ J、$Q=-250$ J。循环总功及净吸热均100 J。
\crosscheck 三段内能变化和为零，三角形面积为100 J；净吸热不等同于正吸热总量950 J。
\sourceproblem{计算7}{244}{250}
\problemtext 1 mol单原子气体矩形循环，a,b,c,d的(V,p)按图为(2,1)、(2,2)、(3,2)、(3,1)，单位分别为10的负3次方立方米、10的5次方Pa，求吸热总量、净功及效率。
\workings ab等容吸热 $Q_{ab}=3(400-200)/2=300$ J；bc等压吸热 $Q_{bc}=5(600-400)/2=500$ J。cd、da均放热，分别为-450及-250 J。因此 $Q_{\rm in}=800$ J，$W_{\rm net}=100$ J，$\eta=100/800=12.5\%$。
\crosscheck 矩形面积100 J，放热总量700 J，第一定律闭合。
\sourceproblem{计算8}{244}{250}
\problemtext 单原子气体初态1 atm、1 L，先等压体积加倍，再等容压强加倍，最后绝热膨胀至初温，求路径、内能变化、吸热、功。
\workings 记初态 $(V_0,p_0)$。B为 $(2V_0,p_0)$，C为 $(2V_0,2p_0)$，温度依次为 $T_0,2T_0,4T_0$。可逆绝热 $TV^{2/3}$ 恒定，D为 $(16V_0,p_0/16)$；p-V图依次画水平线、竖直线、下降绝热线。$\Delta U_{AD}=0$，$Q_{AB}=5p_0V_0/2$，$Q_{BC}=3p_0V_0$，$Q_{CD}=0$。总 $Q=W=(11/2)p_0V_0=557.15$ J。
\crosscheck 末态温度相同而体积不同，并非循环；即使绝热不可逆，末态内能及总功热仍由题给温度和前两段准静态过程确定，末体积则不唯一。
\sourceproblem{计算9}{244}{250}
\problemtext 0.02 kg氦由17摄氏度升至27摄氏度，分别定容、定压、绝热，求内能变化、吸热、外界对气体做功。
\workings $\nu=5$ mol，$\Delta U=3\nu R(10)/2=623.25$ J。定容：$Q=623.25$ J，$W_{\rm on}=0$；定压：$Q=5\nu R(10)/2=1038.75$ J，$W_{\rm on}=-415.50$ J；绝热：$Q=0$，$W_{\rm on}=623.25$ J。
\crosscheck 外界做功与气体对外功符号相反；三种过程内能增量相同。
\sourceproblem{计算10}{244}{250}
\problemtext 1 mol理想气体在400、300 K热源间可逆卡诺循环，高温等温膨胀体积0.001至0.005立方米，求高温吸热、低温放热及净功。
\workings $Q_h=400R\ln5\simeq5349.77$ J；由可逆卡诺关系 $Q_c/Q_h=300/400$，得放热量大小 $Q_c=300R\ln5\simeq4012.33$ J，$W=100R\ln5\simeq1337.44$ J。
\crosscheck 效率25%；低温放热量取正大小，工质有向吸热为负。
\sourceproblem{计算11}{245}{251}
\problemtext 三角循环图中A、B、C的(V,p)分别(1,300)、(3,100)、(1,100)，体积为立方厘米、压强Pa，A温度300 K，求B、C温度及循环净功净热。
\workings $T\propto pV$，故 $T_B=300$ K，$T_C=100$ K。顺时针三角面积 $W=\frac12(300-100)(3-1)10^{-6}=2.0\times10^{-4}$ J，$Q_{\rm net}=W$。
\crosscheck 图示单位是Pa与cm³，没有隐含10的5次方压强因子。
\sourceproblem{计算12}{245}{251}
\problemtext 热容比1.40，A温度300 K，图中A、B、C的(V,p)分别(2,400)、(6,100)、(2,100)，单位cm³、Pa，求B、C温度及各段吸热。
\workings $T_B=225$ K，$T_C=75$ K。定热容下 $\Delta U=\Delta(pV)/(\gamma-1)$。AB直线 $W=1.0\times10^{-3}$ J，$\Delta U=-5.0\times10^{-4}$ J，$Q=5.0\times10^{-4}$ J；BC等压 $W=-4.0\times10^{-4}$ J，$\Delta U=-1.0\times10^{-3}$ J，$Q=-1.4\times10^{-3}$ J；CA等容 $W=0$，$\Delta U=Q=1.5\times10^{-3}$ J。
\crosscheck 净热6.0乘10的负4次方J，等于三角形面积；原题所问各段净热取有向值。
\sourceproblem{计算13}{245}{251}
\problemtext 卡诺机高温源400 K，每循环吸热100 J、放热80 J，求低温源温度和效率。
\workings $T_c=T_hQ_c/Q_h=320$ K，$\eta=1-Q_c/Q_h=20\%$。
\crosscheck 净功20 J。
\sourceproblem{计算14}{245}{251}
\problemtext 热机效率25%，输出功率5 kW，每循环放热8000 J，求吸热及周期。
\workings $Q_h=Q_c/(1-\eta)=10666.67$ J，$W=\eta Q_h=2666.67$ J，$\tau=W/P=0.5333$ s。
\crosscheck 功率是单位时间净功，不是吸热率。
\sourceproblem{计算15}{246}{252}
\problemtext 卡诺机效率40%，500 K热源每循环供热800 J，求放热及低温源温度。
\workings $Q_c=(1-\eta)Q_h=480$ J，$T_c=(1-\eta)T_h=300$ K。
\crosscheck 净功320 J。
\sourceproblem{计算16}{246}{252}
\problemtext 逆斯特林循环，AB、CD分别为高温T1等温压缩、低温T2等温膨胀，BC、DA等容，求制冷系数。
\workings 取 $r=V_2/V_1>1$，低温等温吸热 $Q_c=\nu RT_2\ln r$；净输入功 $W_{\rm in}=\nu R(T_1-T_2)\ln r$。若采用理想完全回热，或题目仅计CD从制冷对象吸收的热量，则 $e=Q_c/W_{\rm in}=T_2/(T_1-T_2)$。
\crosscheck 等容段热量为正负相等的内能变化；实际无回热装置需明确这些热量与哪个热源交换，不能无条件将系统实际制冷系数等同于卡诺值。
\sourceproblem{计算17}{246}{252}
\problemtext 1 mol单原子气体T-V循环，A、B在过原点直线上，V_A=2V_B，BC等容、CA等温，T_C=600 K，求各段热量、净功、效率。
\workings $T_A=600,T_B=300$ K；AB等压压缩，$Q_{AB}=\frac52R(-300)=-6232.5$ J。BC等容升温，$Q_{BC}=\frac32R(300)=3739.5$ J。CA等温膨胀，$Q_{CA}=600R\ln2\simeq3455.30$ J（用题给0.693）。$W_{\rm net}=962.30$ J，$\eta=W_{\rm net}/(Q_{BC}+Q_{CA})\simeq13.38\%$。
\crosscheck 图为T-V而非p-V，不能直接用图形面积表示功；净热与净功相等。
\sourceproblem{计算18}{246}{252}
\problemtext 证明可逆卡诺逆循环制冷系数。
\workings 低温源取热 $Q_c>0$，高温源放热 $Q_h>0$。可逆循环满足 $Q_h/T_1=Q_c/T_2$，能量守恒给 $W_{\rm in}=Q_h-Q_c$。故 $e=Q_c/W_{\rm in}=T_2/(T_1-T_2)$。
\crosscheck 温度采用K，T1大于T2；制冷系数与热泵供热系数相差1。
\sourceproblem{应用1}{247}{253}
\problemtext 瀑布落差979 m，重力功50%转化为水内能，水比热4180 J每千克每K，求温升。
\workings $mc\Delta T=0.5mgh$，故 $\Delta T=0.5(9.8)(979)/4180\simeq1.148$ K。
\crosscheck 质量抵消；题给能量转化比例已包含其他能量去向。
\sourceproblem{应用2}{247}{253}
\problemtext 烟气由900摄氏度降至200摄氏度，给密度1.29 kg每立方米、摩尔质量0.0289 kg/mol及定压摩尔热容7R/2，求单位体积放热。
\workings 按所给密度对应的体积计量，1立方米含 $\nu=1.29/0.0289$ mol，定压放热大小 $\nu(7R/2)(700)\simeq9.09\times10^5$ J。
\crosscheck 密度与体积必须对应同一参考状态；若问实际900摄氏度入口立方米或200摄氏度出口立方米，而1.29为标准态空气密度，则需按状态方程重算。
\sourceproblem{应用3}{247}{253}
\problemtext 2 mol氦气初温300 K、20 L，等压膨胀体积加倍，再绝热膨胀至初温，求路径、内能、吸热、功。
\workings 初态 $p_A=2R(300)/0.020=249300$ Pa，B体积0.040立方米、温度600 K、压强同A。若可逆绝热，C体积 $0.040(600/300)^{3/2}\simeq0.1131$ 立方米，压强 $2R(300)/V_C\simeq4.407\times10^4$ Pa。画水平AB及下降绝热线BC。$\Delta U=0$，$Q=2(5R/2)(300)=12465$ J，$W=12465$ J。
\crosscheck 等压功4986 J，绝热功7479 J；不可逆绝热时总功热仍相同，但C体积不能按绝热状态关系确定。
\sourceproblem{应用4}{247}{253}
\problemtext 1 mol刚性氧气初温300 K、1 atm，分别等温或绝热膨胀到三倍体积，求对外功。
\workings 准静态等温 $W_T=RT\ln3\simeq2738.81$ J。可逆绝热 $\gamma=7/5$，$T_f=T_i3^{-2/5}$，$W_S=\frac52RT_i(1-3^{-2/5})\simeq2216.31$ J。
\crosscheck 两者并非同末态；仅指定不可逆绝热的末体积不能唯一确定功，主解按教材可逆路径。
\sourceproblem{应用5}{248}{254}
\problemtext 原书以高压气体推动物体作绝热膨胀为背景，要求分析气体功及物体动能变化。
\workings 一般准静态绝热过程满足 $pV^\gamma=C$。对非武器化的封闭活塞热力学模型，气体功为 $W=\int_{V_i}^{V_f}p\,dV=(p_iV_i-p_fV_f)/(\gamma-1)$；若仅存在理想能量转移，则活塞动能变化由净外力功确定。
\crosscheck 本题保留原书编号与通用热力学推导，不给出针对原武器装置的性能数值；原题数值部分列为未完成，不计入完整数值解析。
\sourceproblem{应用6}{248}{254}
\problemtext 地热热机高、低温源227及27摄氏度，每小时吸热1.8乘10的11次方J，求最大功率。
\workings $\eta_{\max}=1-300/500=0.4$，吸热率 $1.8\times10^{11}/3600=5.0\times10^7$ W，故最大输出功率 $2.0\times10^7$ W，即20 MW。
\crosscheck 卡诺上限，实际发电还受转换及辅机损耗限制。
\sourceproblem{应用7}{248}{254}
\problemtext 室外37摄氏度、室内17摄氏度，每日漏入热量2.51乘10的8次方J，实际制冷系数为卡诺值60%，求每日耗电。
\workings $e_C=290/(310-290)=14.5$，$e=0.6e_C=8.7$。$W=2.51\times10^8/8.7\simeq2.885\times10^7$ J，换算约8.014 kWh。
\crosscheck 按题设恒温及仅补偿给定热负荷，不含其他设备耗能。
